% ============================================================
% Appendix F — Governance and Law (Phase-7)
% ============================================================

%\subsection{Governance, Law, and Epistemic Oversight}

\subsubsection{Position in the Architecture}

Previously we establish:

\begin{itemize}
\item individual pipeline correctness  ,
\item dynamic degradation and hallucination inevitability  ,
\item collective epistemic dynamics and population effects  ,
\item formal impossibility results  ,
\item institutional stabilization mechanisms .
\end{itemize}

Next we study \emph{governance and law} as higher-order structures that
operate \emph{on institutions themselves}, rather than on pipelines or agents.

\subsubsection{Governance as a Meta-Institution}
\label{sec:appendix-F1}

Let $J$ and $M(J)$ be as in Appendix~E, and let an \emph{institutional aggregator family} be a collection
\[
  (I_P)_P, \qquad I_P : M(J) \to J \cup \{\bot\},
\]
assigning to each FRFP pipeline $P$ an institutional aggregator $I_P$ in the sense of Definition~B.87.

\begin{definition}[Category of institutions]\label{def:F1-Inst}
We define the category $\Inst$ as follows.
\begin{itemize}
  \item Objects are institutional aggregator families $(I_P)_P$.
  \item A morphism $\Phi : (I_P)_P \to (I'_P)_P$ is a natural transformation whose components
  \[
    \Phi_P : M(J) \to M(J)
  \]
  are monotone with respect to refinement of multisets and preserve unanimity: whenever all elements of $S \in M(J)$ agree, so do the elements of $\Phi_P(S)$.
\end{itemize}
Intuitively, a morphism in $\Inst$ is a structure-preserving transformation between institutional layers that does not introduce disagreement where there was none.
\end{definition}

\begin{definition}[Governance operator]\label{def:F1-governance}
A \emph{governance operator} is an endofunctor
\[
  G : \Inst \to \Inst
\]
together with a choice of observable
\[
  \Omega : \mathsf{Pop} \to O
\]
from the category of epistemic populations (Appendix~C) to some space of observables $O$, such that for each time step $t$,
\[
  (I^{(t+1)}_P)_P = G(I^{(t)}_P)_P
\]
is allowed to depend only on the previous institutional layer $(I^{(t)}_P)_P$ and on the history of observables $(\Omega(P_s))_{s \le t}$ generated by the collective dynamics of Appendix~C.
\end{definition}


\subsubsection{Link to Collective Epistemic Dynamics}

Appendix~C models collective epistemic evolution as a dynamical system over
populations:
\[
(\mathbb{P}_t, \Sem_{\mathbb{P}_t})_{t \ge 0}.
\]

Governance operators are defined on this level.

\begin{definition}[Governance observable]
A governance observable is any functional
\[
\Omega : (\mathbb{P}, \Sem_{\mathbb{P}}) \longrightarrow \mathbb{R}
\]
that measures population-level epistemic properties (e.g.\ disagreement
rates, confidence inflation, drift velocity).
\end{definition}

\begin{definition}[Governed evolution]
A governed epistemic system evolves by:
\[
\mathcal{I}_{t+1} = \mathcal{G}(\mathcal{I}_t, \Omega_t),
\qquad
(\mathbb{P}_{t+1}, \Sem_{\mathbb{P}_{t+1}}) = \Phi(\mathcal{I}_{t+1}),
\]
where $\Phi$ is the collective dynamics functor from Appendix~C.
\end{definition}

Thus governance is \emph{feedback on institutions}, not on beliefs.

\subsubsection{Law as a Constraint on Governance}

Governance itself is subject to failure modes: capture, overfitting,
short-term optimization, or erosion of tacit authority.

\begin{definition}[Legal constraint]
A legal constraint is a predicate
\[
\mathcal{L}(\mathcal{I}) \in \{\text{admissible},\text{inadmissible}\}
\]
on institutional forms.
\end{definition}

\begin{definition}[Lawful governance]
A governance operator $\mathcal{G}$ is lawful if:
\[
\mathcal{L}(\mathcal{I}) = \text{admissible}
\;\Rightarrow\;
\mathcal{L}(\mathcal{G}(\mathcal{I})) = \text{admissible}.
\]
\end{definition}

Law therefore restricts the allowable transformations of institutions, not
the content of scientific claims.

\subsubsection{No Fully Automated Governance}

We now state the core limit result.
\begin{definition}[Explicit-only governance]\label{def:F-explicit-only-governance}
A governance operator
\[
  G : \Inst \to \Inst
\]
is \emph{explicit-only} if, for every time $t$ and institutional layer $(I^{(t)}_P)_P$, there exists a family of explicit-only mechanisms
\[
  G^{\mathrm{exp}}_{P,t} : E^{k_{P,t}} \to J \cup \{\bot\}
\]
in the sense of Appendix~D such that the updated institutional aggregators satisfy
\[
  I^{(t+1)}_P(S)
  \;=\;
  G^{\mathrm{exp}}_{P,t}\bigl(e^{(t)}_1, \dots, e^{(t)}_{k_{P,t}}\bigr)
\]
for every multiset $S \in M(J)$, where $e^{(t)}_1, \dots, e^{(t)}_{k_{P,t}} \in E$ are the explicit artifacts produced by the admissible runs of $P$ up to time $t$. In particular, $G$ has no access to any tacit states or tacit correctness outcomes, only to explicit execution traces and their explicit summaries.
\end{definition}
\begin{definition}[Explicit-only governance]\label{def:F-explicit-only-governance}
A governance operator
\[
  G : \Inst \to \Inst
\]
is \emph{explicit-only} if, for every time $t$ and institutional layer $(I^{(t)}_P)_P$, there exists a family of explicit-only mechanisms
\[
  G^{\mathrm{exp}}_{P,t} : E^{k_{P,t}} \to \bigl(M(J) \to J \cup \{\bot\}\bigr)
\]
in the sense of Appendix~D such that the updated institutional aggregators satisfy
\[
  I^{(t+1)}_P(S)
  \;=\;
  G^{\mathrm{exp}}_{P,t}\bigl(e^{(t)}_1, \dots, e^{(t)}_{k_{P,t}}\bigr)(S)
\]
for every multiset $S \in M(J)$, where $e^{(t)}_1, \dots, e^{(t)}_{k_{P,t}} \in E$ are the explicit artifacts produced by the admissible runs of $P$ up to time $t$. In particular, $G$ has no access to any tacit states or tacit correctness outcomes, only to explicit execution traces and their explicit summaries.
\end{definition}

\begin{theorem}[Non-automation of governance]
No governance operator $\mathcal{G}$ that depends solely on explicit artifacts
or explicitly computable observables can preserve epistemic correctness over
time.
\end{theorem}

\begin{proof}[Proof sketch]
By Appendix~B, explicit-only observables cannot detect tacit degradation.
By Appendix~D, correctness depends on $\gamma(\sem{r})$, which is tacit.
Any governance operator computable from explicit traces alone therefore
cannot reliably detect institutional epistemic failure modes.
\end{proof}

\subsubsection{Human Oversight as a Structural Requirement}

\begin{corollary}[Human oversight necessity]
Any lawful governance operator must incorporate human tacit judgment at the
meta-institutional level.
\end{corollary}

This does not imply that humans must make all decisions, but that \emph{some
decisions cannot be delegated}.

\subsubsection{Regulatory Scope}

Within FRFP, regulation can only legitimately target:

\begin{itemize}
\item admissible institutional forms,
\item governance feedback mechanisms,
\item transparency of explicit pipelines,
\item preservation of tacit authority.
\end{itemize}

It cannot:
\begin{itemize}
\item certify correctness automatically,
\item replace tacit evaluation with metrics,
\item eliminate hallucination in principle.
\end{itemize}

\subsubsection{Final Synthesis}

FRFP yields a stratified epistemic architecture:

\begin{center}
\begin{tabular}{c}
\textbf{Pipelines} — compute explicit artifacts \\
\textbf{Agents} — perform tacit judgment \\
\textbf{Institutions} — stabilize populations \\
\textbf{Governance} — adapt institutions \\
\textbf{Law} — constrain governance
\end{tabular}
\end{center}

Each layer is necessary; none subsumes the one above it.

\medskip

\noindent
\emph{There is no fully automated science, no fully automated governance of
science, and no lawful system that claims otherwise.}

% End Appendix F
